
\begin{frame}{Aggregation, Learning, Evaluation and Deployment Bias}

    \begin{itemize}
        \setlength\itemsep{0.7em}
        \item \textbf{Aggregation bias.} One model for populations that genuinely behave
              differently. Suresh and Guttag point to medicine: patients of different sexes with
              similar underlying conditions can present and progress differently, so one model is
              optimal for neither.

        \item \textbf{Learning bias.} The modelling choice itself amplifies a disparity.
              Optimising average loss lets a large majority group dominate the gradient;
              aggressive compression or regularisation tends to cost the tails most.

        \item \textbf{Evaluation bias.} The benchmark does not represent the use population,
              so the reported number is not the number you will get. This is exactly the
              Gender Shades finding about existing face benchmarks.

        \item \textbf{Deployment bias.} The model is used differently from how it was framed.
              A risk score built to \emph{rank} people for services becomes a
              \emph{decision} about their liberty; a triage aid becomes an autopilot.
    \end{itemize}

    \vspace{0.4em}
    \centering
    \textit{Four of the seven sources arise \emph{after} the data has been collected.}
\end{frame}
